% !TEX program = xelatex
% ============================================================
%  毕业答辩讲稿（独立成册）
%  编译：xelatex script.tex（跑两次）
%  节与 slides.tex 的 20 张幻灯片一一对应
% ============================================================
\documentclass[a4paper,11pt]{ctexart}
\usepackage[margin=2.2cm]{geometry}
\usepackage{enumitem}
\setlist{noitemsep, topsep=0.2em, leftmargin=1.5em}

% ---------- 共享 preamble ----------
% !TEX program = xelatex
% ---------- 共享 preamble ----------
% slides 与 script 共用。仅拉 _figpreamble.tex 的颜色宏 + tikz + pgfplots + subcaption。
% 不 \input 毕业论文/preamble.tex（避免 ctexbook 与 ctexbeamer/ctexart 冲突）。

% ---------- 字体 ----------
\def\ProjectFontDir{../fonts}
% Project-local font setup for XeLaTeX.
% Callers may set \ProjectFontDir before \input'ing this file.
\providecommand{\ProjectFontDir}{../fonts}

\setmainfont[
  Path = \ProjectFontDir/tex-gyre/,
  Extension = .otf,
  UprightFont = *-regular,
  BoldFont = *-bold,
  ItalicFont = *-italic,
  BoldItalicFont = *-bolditalic,
  Ligatures = TeX
]{texgyretermes}

\setCJKmainfont[
  Path = \ProjectFontDir/fandol/,
  Extension = .otf,
  BoldFont = FandolSong-Bold,
  ItalicFont = FandolKai-Regular
]{FandolSong-Regular}
\setCJKsansfont[
  Path = \ProjectFontDir/fandol/,
  Extension = .otf,
  BoldFont = FandolHei-Bold,
  ItalicFont = FandolKai-Regular
]{FandolHei-Regular}
\setCJKmonofont[
  Path = \ProjectFontDir/fandol/,
  Extension = .otf,
  ItalicFont = FandolKai-Regular
]{FandolFang-Regular}

\setCJKfamilyfont{zhsong}[
  Path = \ProjectFontDir/fandol/,
  Extension = .otf,
  BoldFont = FandolSong-Bold
]{FandolSong-Regular}
\setCJKfamilyfont{zhsongb}[
  Path = \ProjectFontDir/fandol/,
  Extension = .otf
]{FandolSong-Bold}
\setCJKfamilyfont{zhhei}[
  Path = \ProjectFontDir/fandol/,
  Extension = .otf,
  BoldFont = FandolHei-Bold,
  ItalicFont = FandolKai-Regular
]{FandolHei-Regular}
\setCJKfamilyfont{zhheib}[
  Path = \ProjectFontDir/fandol/,
  Extension = .otf
]{FandolHei-Bold}
\setCJKfamilyfont{zhfs}[
  Path = \ProjectFontDir/fandol/,
  Extension = .otf
]{FandolFang-Regular}
\setCJKfamilyfont{zhkai}[
  Path = \ProjectFontDir/fandol/,
  Extension = .otf
]{FandolKai-Regular}

\ExplSyntaxOn
\cs_if_exist:NT \ctex_punct_set:n
  {
    \ctex_punct_set:n { fandol }
    \ctex_punct_map_family:nn   { \CJKrmdefault         } { zhsong  }
    \ctex_punct_map_family:nn   { \CJKsfdefault         } { zhhei   }
    \ctex_punct_map_family:nn   { \CJKttdefault         } { zhfs    }
    \ctex_punct_map_bfseries:nn { \CJKrmdefault, zhsong } { zhsongb }
    \ctex_punct_map_bfseries:nn { \CJKsfdefault, zhhei  } { zhheib  }
    \ctex_punct_map_itshape:nn  { \CJKrmdefault         } { zhkai   }
  }
\ExplSyntaxOff

\providecommand{\songti}{}
\providecommand{\heiti}{}
\providecommand{\fangsong}{}
\providecommand{\kaishu}{}
\renewcommand{\songti}{\CJKfamily{zhsong}}
\renewcommand{\heiti}{\CJKfamily{zhhei}}
\renewcommand{\fangsong}{\CJKfamily{zhfs}}
\renewcommand{\kaishu}{\CJKfamily{zhkai}}


% ---------- 宏包 ----------
\usepackage{amsmath,amssymb}
\usepackage{booktabs}
\usepackage{array}
\usepackage{multirow}
\usepackage{tikz}
\usetikzlibrary{shapes.geometric, arrows.meta, positioning, fit, backgrounds, calc, patterns, decorations.pathreplacing}
\usetikzlibrary{external}
\tikzexternalize
\usepackage{pgfplots}
\pgfplotsset{compat=1.18}
\usepgfplotslibrary{fillbetween}
\pgfplotsset{table/search path={../图片/, ../图片/data/, ../}}
\usepackage{xcolor}
\usepackage{colortbl}
\usepackage{caption}
\usepackage{subcaption}
\usepackage{adjustbox}
\usepackage[absolute,overlay]{textpos}
\usepackage[skins]{tcolorbox}

% ---------- 统一调色板（与 图片/_figpreamble.tex 同步） ----------
\definecolor{deepblue}{RGB}{81, 132, 178}
\definecolor{lightblue}{RGB}{170, 212, 248}
\definecolor{skyblue}{RGB}{154, 220, 255}
\definecolor{inputyellow}{RGB}{255, 248, 154}
\definecolor{outputpink}{RGB}{255, 178, 166}
\definecolor{improvedpink}{RGB}{255, 138, 174}
\definecolor{lightpink}{RGB}{241, 167, 181}
\definecolor{deeppink}{RGB}{213, 82, 118}
\definecolor{bglight}{RGB}{242, 245, 250}
\definecolor{lightgreen}{RGB}{180, 230, 180}
\definecolor{lightorange}{RGB}{255, 210, 140}
\definecolor{dangerred}{RGB}{230, 100, 100}
\definecolor{vibrantgreen}{RGB}{46, 175, 80}
\definecolor{vibrantorange}{RGB}{240, 130, 30}
\definecolor{vibrantred}{RGB}{220, 50, 50}

% ---------- 答辩页浮动信息卡 ----------
\tcbset{
  slideMetricCard/.style={
    enhanced,
    colback=white,
    colframe=deepblue!70!black,
    coltitle=white,
    colbacktitle=deepblue!85!black,
    fonttitle=\bfseries\scriptsize,
    boxrule=0.5pt,
    arc=2pt,
    left=4pt,
    right=4pt,
    top=3pt,
    bottom=0pt,
    toptitle=1pt,
    bottomtitle=1pt,
    titlerule=0pt,
    drop fuzzy shadow=black!18
  }
}

% ---------- 统一 tikz 样式（与 图片/_figpreamble.tex 同步） ----------
\tikzset{
  mod/.style            = {rectangle, rounded corners=2pt, draw=deepblue, fill=lightblue, minimum height=0.9cm, align=center, font=\small},
  modimp/.style         = {rectangle, rounded corners=2pt, draw=deeppink, fill=improvedpink, minimum height=0.9cm, align=center, font=\small},
  modin/.style          = {rectangle, draw=deepblue, fill=inputyellow, minimum height=0.9cm, align=center, font=\small},
  modout/.style         = {rectangle, draw=deeppink, fill=outputpink, minimum height=0.9cm, align=center, font=\small},
  axx/.style            = {->, thick, black!80},
  flow/.style           = {-Stealth, thick, deepblue},
  flowimp/.style        = {-Stealth, thick, deeppink},
  obsstatic/.style      = {fill=dangerred!70, draw=dangerred, thick},
  obsdyn/.style         = {fill=lightorange, draw=dangerred, thick},
  obsbuf/.style         = {draw=dangerred!70, dashed, thick},
  bandok/.style         = {fill=lightgreen, opacity=0.55},
  bandbad/.style        = {fill=dangerred!30, opacity=0.55},
  egopt/.style          = {fill=deepblue, draw=deepblue},
  reflinestyle/.style   = {deepblue, very thick},
}
% ---------- 论文实验宏（与 毕业论文/ 同步） ----------
\InputIfFileExists{../毕业论文/_experiment_metrics.tex}{}{}
\InputIfFileExists{../毕业论文/_experiment_context.tex}{}{}
\InputIfFileExists{../毕业论文/_ablation_macros.tex}{}{}
\InputIfFileExists{../毕业论文/_sensitivity_macros.tex}{}{}

% ---------- 算法常量宏（slides 层级，非自动生成） ----------
\newcommand{\GroupKmax}{15}
\newcommand{\GroupCThreeK}{18}
\newcommand{\GroupCThreeCombCount}{26}
\newcommand{\ComplexityOldBase}{$O(2^k)$}
\newcommand{\ComplexityNewBase}{$O(k\log k)$}


\usepackage{hyperref}
\hypersetup{colorlinks=true, linkcolor=deepblue, urlcolor=deepblue,
            pdftitle={毕业答辩讲稿\quad{}MIKU 算法},
            pdfauthor={廖嘉旺}}

% 节标题样式（不用冒号）
\usepackage{titlesec}
\newcommand{\sectionbreak}{\clearpage}
\titleformat{\section}[block]
  {\normalfont\Large\bfseries\color{deepblue}}{Slide \thesection}{1em}{}
\titlespacing*{\section}{0pt}{1.2em}{0.6em}

\setcounter{secnumdepth}{1}

\title{\textbf{毕业答辩讲稿}\\[0.4em]\large 基于Apollo的多障碍物场景下路径规划算法实现}
\author{廖嘉旺\quad{}指导教师\quad{}孙成娇}
\date{2026年5月}

\begin{document}

\maketitle

\vspace{0.5em}
\noindent\textbf{讲稿使用说明}\quad{}本讲稿与 \texttt{slides.pdf} 的 19 张幻灯片逐张对应（Apollo 现状节已合并至 MIKU 总体框架，作为追问储备保留在本稿）。每节给出该 Slide 的讲述时间窗口、口语要点、过渡句与可能追问。节标题 ``Slide XX'' 中 XX 为幻灯片编号。累计时间以封面起算，最后一页收在 10 分钟之内，保留约 30 秒缓冲。

\section{封面}
\textbf{时间}\quad{}0:00 -- 0:20（累计 00:20）\\
\textbf{要点}
\begin{itemize}
  \item 各位老师好，我叫廖嘉旺，来自计科 2211 班，指导教师孙成娇老师。
  \item 今天汇报的题目是\emph{基于 Apollo 的多障碍物场景下路径规划算法实现}，我把本文提出的算法统一命名为 MIKU。
  \item 整个汇报大约 10 分钟，分三部分，先讲问题与目标，再讲三点改进，最后讲实验与结论。
\end{itemize}
\textbf{过渡}\quad{}下面从研究背景讲起。\\
\textbf{可能追问}\quad{}MIKU 是什么缩写。应对，Multi-obstacle Interaction-density Kinematic-prediction Unified Planner，意思是多障碍物交互运动学统一规划算法。

\section{研究背景}
\textbf{时间}\quad{}0:20 -- 1:05（累计 01:05）\\
\textbf{要点}
\begin{itemize}
  \item 自动驾驶从感知到控制中间这一段，路径规划直接决定整车的安全与舒适。
  \item 多障碍物密集场景有三个典型难点，第一可行域收缩，几个禁行区一叠加通行空间就压到零，第二决策耦合，避让方向之间会相互制约，第三时间预算紧，规划周期不到 100 毫秒。
  \item 本文把研究基座定在百度 Apollo 11.0 的 PublicRoad Planner，具体切入点是路径边界决策模块 PathBoundsDecider。
\end{itemize}
\textbf{过渡}\quad{}那 Apollo 现有做法到底哪里次优，看下一张。\\
\textbf{可能追问}\quad{}为什么不用学习方法。应对，学习方法在安全保证与形式化验证上仍有距离，本文定位是在工业级凸优化框架内给增量改进，与学习方法是互补关系。

\section{问题陈述}
\textbf{时间}\quad{}1:05 -- 1:50（累计 01:50）\\
\textbf{要点}
\begin{itemize}
  \item 左图是 Baseline 贪心决策流程，我通过源码走查得到三个观察。
  \item 第一，安全裕度一刀切，\texttt{GetBufferBetweenADCCenterAndEdge} 对所有障碍物返回半车宽加 0.4 米，路锥、花坛、违停车一视同仁。
  \item 第二，独立决策矛盾，\texttt{UpdatePathBoundaryBySLPolygon} 对每个障碍物独立决定 LEFT\_NUDGE 或 RIGHT\_NUDGE，容易把左右两侧方向拧反。
  \item 第三，动态障碍物排除，\texttt{IsStatic} 过滤让行人在路径阶段完全不可见。
\end{itemize}
\textbf{过渡}\quad{}三个症状分别对应安全距离、决策粒度、时空协调，这就是三条改进维度。\\
\textbf{可能追问}\quad{}Apollo 的 LaneBorrowPath 能不能兜底。应对，LaneBorrowPath 要等十几个周期才触发，不是及时的解决方案，且借道本身会打扰邻道，本文目标是在本车道内先找最优通行带。

\section{研究目标与贡献}
\textbf{时间}\quad{}1:50 -- 2:20（累计 02:20）\\
\textbf{要点}
\begin{itemize}
  \item 对应三条维度给出三点改进，多因子威胁度、组级最优通行带、时空可通行走廊。
  \item 记忆点是\textbf{一组符号、三项改进、一步 PathBoundsDecider}。三项改进共用同一套 $(u_i, v_i, \delta, \tau)$ 符号，不是三块补丁；源码层只动 Apollo 的 \texttt{PathBoundsDecider} 一步，PJPO、PJSO 与 Control 链路零改动。
  \item 组级最优把 $O(2^k)$ 的暴力搜索压到 $O(k\log k)$ 的精确求解。
\end{itemize}
\textbf{过渡}\quad{}在讲改进之前先交代 Apollo 现状，给大家定位本文改动的范围。\\
\textbf{可能追问}\quad{}为什么要证明最优性。应对，答辩场上要硬对比，组级最优需要引理和定理兜住，否则就是启发式策略。

\section{Apollo Planning 模块现状}
\textbf{时间}\quad{}合并至上节，作为 MIKU 总体框架的铺垫。若评委追问再展开。\\
\textbf{要点}
\begin{itemize}
  \item Apollo 的 Planning 采用 Scenario-Stage-Task 三级管理架构，LaneFollow、LaneChange、ParkAndGo 等场景都在这套架构里。
  \item 路径阶段由分段 Jerk 路径优化器 PJPO 承担，前置的 PathBoundsDecider 把障碍物决策转成 QP 的左右边界。
  \item 本文所有改动都限定在路径边界构建这一步，上游感知、下游 PJPO 与控制模块都不动。
\end{itemize}
\textbf{备注}\quad{}本节已从 slides 中裁掉，只保留讲稿内容作为追问储备。\\
\textbf{可能追问}\quad{}为什么不动速度 QP。应对，本文的第三点改进确实会把时空走廊注入速度 QP 的上界，但速度 QP 自身的求解器 PJSO 不改。

\section{MIKU 总体框架}
\textbf{时间}\quad{}2:20 -- 3:40（累计 03:40）\\
\textbf{要点}
\begin{itemize}
  \item 整体流水三阶段，与论文 5/6/7 章一一对应。
  \begin{itemize}
    \item \textbf{裕度级}（第5章）——多因子威胁度 $\Theta_i$ 映射到分级 $d_{buf,i}$，得差异化 $\delta_i$。
    \item \textbf{组级}（第6章）——扫描线分组 $\mathcal{G}_j$，组内按 $u_i$ 升序取最大间隙得 $\mathbf{d}^*$。
    \item \textbf{时空级}（第7章）——到达时间 $\tau(s)$ 投影 $O_i(\tau(s))$，并把组级通行带的时间有效性向速度 QP 反投影为 ST 走廊。
  \end{itemize}
  \item 输出仍然是 Apollo 标准的 \texttt{PathBoundary} 与 \texttt{STDrivableBoundary}，下游模块零改动。
\end{itemize}
\textbf{过渡}\quad{}下面分别展开三点改进。\\
\textbf{可能追问}\quad{}这个流水会不会超预算。应对，扫描线分组 $O(n\log n)$ 加逐离散点 $O(K)$，实测 QP 总耗时远低于 Baseline，第 16 张会给出数字。

\section{创新一\quad{}多因子威胁度}
\textbf{时间}\quad{}3:40 -- 4:30（累计 04:30）\\
\textbf{要点}
\begin{itemize}
  \item 把一刀切的额外缓冲 \DeltaBaseline\,米换成分级 $d_{buf,i} \in [\DeltaIntervalLow, \DeltaIntervalHigh]$\,米（总裕度 $\delta_i = W_{ego}/2 + d_{buf,i}$），由 5 个因子加权决定。
  \item TTC 反映碰撞紧迫度，横向重叠率反映是否真的挡路，相对速度 Sigmoid 映射接近速度大小，障碍物类型先验区分锥桶与违停车，交互密度体现是否处于人群中。
  \item 5 项通过固定权重加权归一化为 $\Theta_i\!\in\![0,1]$，$d_{buf,i} = d_{buf,\min} + (d_{buf,\max}-d_{buf,\min})\cdot \Theta_i$，锥桶 $\Theta_i\!\approx\!0.2$，行人 $\Theta_i\!\approx\!0.9$。
\end{itemize}
\textbf{过渡}\quad{}这个方法在窄路锥桶场景下效果最直接。\\
\textbf{可能追问}\quad{}权重怎么定。应对，固定权重来自初始参数调优没有在线学习，灵敏度分析显示 $\pm\SensPerturbPct\%$ 扰动下最坏相对偏差仅 \SensWorstRelDev\%。

\section{创新一\quad{}实验验证（窄路密集锥桶）}
\textbf{时间}\quad{}4:30 -- 5:10（累计 05:10）\\
\textbf{要点}
\begin{itemize}
  \item P3 窄路锥桶场景，左右各 3 个水马/锥桶（共 6 个）沿车道两侧布置成通行带。
  \item Baseline 统一 0.30 米额外缓冲，左右净间距不足车宽，被卡在 $s\!=\!\MThreeBaselineSEnd$\,m 处停车 \MThreeBaselineStopDur\,秒，最终 TArrive 是 $\infty$，没走完。
  \item MIKU 差异化裕度把锥桶的 $d_{buf,i}$ 压到 0.10 米，打开一条连续通行带，14 秒完成通行，通行效率从 \MThreeBaselineEfficiency\,提升到 \MThreeMikuEfficiency。
  \item 平顺度提升更直观，最大 Jerk 从 \MThreeBaselineMaxAbsJerk\,m/s\textsuperscript{3} 降到 \MThreeMikuMaxAbsJerk\,m/s\textsuperscript{3}，接近 5.7 倍。
\end{itemize}
\textbf{过渡}\quad{}这张图只改了裕度参数，下面讲决策粒度。\\
\textbf{可能追问}\quad{}0.10 米是不是过激。应对，0.10 米是差异化裕度下界，对应静止低威胁物；Apollo 源码 \texttt{obstacle\_lat\_buffer} 默认 0.4 米，本文仿真按论文 chapter 8 口径取 0.30 米，两者均在可接受安全范围内。

\section{创新二\quad{}组级最优通行带}
\textbf{时间}\quad{}5:10 -- 6:00（累计 06:00）\\
\textbf{要点}
\begin{itemize}
  \item 贪心容易拧反的根源是每个障碍物独立决策，本文把同一组在纵向上重叠的障碍物看作一个连通分量。
  \item 引理证明，最优分割在按横向排序后一定是连续的，不会交错。
  \item 定理进一步说，组内按 $u_i$ 升序（动态 $\delta_i$ 使 $u_i=l_i^{-}-\delta_i$ 的排序未必等价 $l_i^{-}$）得到 $k+1$ 个候选间隙，最大那一个间隙就是最优解。
  \item 复杂度由 $O(2^k)$ 指数级降至 $O(k\log k)$ 对数线性。以 P4 密集施工为例，组规模 $k\!=\!\GroupCThreeK$，原本 $2^{\GroupCThreeK}\!\approx\!\GroupCThreeCombCount$ 万次组合判断，MIKU 按排序一次给出最优解。
\end{itemize}
\textbf{过渡}\quad{}下一张给出最典型的行人加停车场景。\\
\textbf{可能追问}\quad{}连通分量太大怎么办。应对，论文在降级机制里设置 $k_{\max}\!=\!15$ 阈值，超过后回退到 Baseline 的逐障碍物策略。

\section{创新二\quad{}实验验证（行人加违停车）}
\textbf{时间}\quad{}6:00 -- 6:40（累计 06:40）\\
\textbf{要点}
\begin{itemize}
  \item P2 行人加违停车场景，左侧违停车，右侧一位横向行走的行人。
  \item Baseline 独立决策给左侧违停分配 RIGHT\_NUDGE，给右侧行人分配 LEFT\_NUDGE，方向对冲导致 $l^+<l^-$，路径 BLOCKED，到达时间 $\infty$。
  \item MIKU 组级决策把两者看成一个连通分量，中间间隙最宽，自车从两者之间穿过，\MTwoMikuTArrive\,秒到位，通行效率从 \MTwoBaselineEfficiency\,提升到 \MTwoMikuEfficiency。
  \item 本场景 QP 耗时差距最显著，Baseline 因路径阻塞触发多轮速度 QP 回退，总耗时 \MTwoBaselineQpTotal\,ms；MIKU 一次收敛只要 \MTwoMikuQpTotal\,ms，加速 20.6 倍。
\end{itemize}
\textbf{过渡}\quad{}行人刚才按静态处理，下面讲真正的动态场景。\\
\textbf{可能追问}\quad{}为什么这里把行人当成静态。应对，SL 平面上行人横向位置在一个规划周期内近似静态，动态特性在下一点的时空走廊里展开。

\section{创新三\quad{}时空可通行走廊}
\textbf{时间}\quad{}6:40 -- 7:30（累计 07:30）\\
\textbf{要点}
\begin{itemize}
  \item 引入到达时间函数 $\tau(s)$，表示自车预估到达位置 $s$ 的时刻，把规划维度从 SL 二维升格为 ST 时空。$\tau(s)$ 由上周期速度曲线查表获得，失效时退回二阶运动学模型。
  \item 两路约束融合。其一为时变 SL 投影 $O_i(\tau(s))$，动态障碍物的预测轨迹参与路径边界构建，Apollo 原来的 \texttt{IsStatic} 过滤在路径阶段被废弃。
  \item 其二为速度 QP 反投影，组级通行带的时间有效性反投影为速度 QP 的 $s$ 上界 $s_j^{ub}$，与原 SBD 边界取交集后收紧。
  \item 量化效益可以用 P1 单行人看，Baseline TArrive 是 $\infty$，MIKU \MOneMikuTArrive\,秒；P4 密集施工 Baseline 效率 \MFourBaselineEfficiency，MIKU \MFourMikuEfficiency。
\end{itemize}
\textbf{过渡}\quad{}下一张在 ST 图上看效果。\\
\textbf{可能追问}\quad{}$\tau(s)$ 怎么来。应对，单次前向估计，由上一帧速度曲线反推或运动学模型给出，不是迭代优化，论文承认这是局限。

\section{创新三\quad{}实验验证（单行人横穿）}
\textbf{时间}\quad{}7:30 -- 8:10（累计 08:10）\\
\textbf{要点}
\begin{itemize}
  \item P1 单行人横穿场景，行人以横向速度 $v_{lat}$ 从车道中线向左横穿。
  \item Baseline 路径阶段因 \texttt{IsStatic} 过滤不避让行人，到了速度阶段在 ST 图上反复试探，触发多轮 fallback，TArrive 是 $\infty$。
  \item MIKU 由 $\tau(s)$ 预判行人在自车到达时已横出车道，速度 QP 一次收敛，\MOneMikuTArrive\,秒到位，通行效率从 \MOneBaselineEfficiency\,提升到 \MOneMikuEfficiency。
  \item 最大 Jerk 从 \MOneBaselineMaxAbsJerk\,m/s\textsuperscript{3} 降到 \MOneMikuMaxAbsJerk\,m/s\textsuperscript{3}，降 5.1 倍。
\end{itemize}
\textbf{过渡}\quad{}这是三点改进各自的验证，下面看综合实验。\\
\textbf{可能追问}\quad{}预测错了怎么办。应对，论文第三级降级是预测轨迹不可用时退化为静态处理，保证不劣于 Baseline。

\section{综合实验设置}
\textbf{时间}\quad{}8:10 -- 8:40（累计 08:40）\\
\textbf{要点}
\begin{itemize}
  \item 场景矩阵总共 \MScenarioTotal\,组，其中压力场景 \MPressureScenarioTotal\,组依次是 P1 单行人横穿、P2 行人加停车、P3 窄路锥桶、P4 密集施工；可比场景 \MComparableScenarioTotal\,组 C1--C4 对应同一几何但调整到双方均可达。
  \item 消融给出 \AblationVariantTotal\,组变体，M0 是 Apollo Baseline，M5 是 MIKU 全量，M1--M4 是依次关闭 C1/C2--C3/C4/C5 的减法消融（leave-one-out），用于定位每个组件的边际贡献。
  \item 评分体系 \ScoreDimensionCount\,维 \ScoreMetricCountTotal\,指标，权重分别是效率 \ScoreWeightEfficiency、平滑 \ScoreWeightSmoothness、鲁棒 \ScoreWeightRobustness、计算 \ScoreWeightCompute。
  \item 灵敏度分析是威胁度权重 $\pm\SensPerturbPct\%$ 扰动 \SensTrialCount\,次蒙特卡洛，用来判断策略对超参的依赖。
  \item 所有数据用 Python 复现 apollo\_pipeline，OSQP 实测 QP 耗时，CSV 和 JSON 入库可一键重跑。
\end{itemize}
\textbf{过渡}\quad{}先看通过率。\\
\textbf{可能追问}\quad{}为什么不在真车上跑。应对，验证限于 Python 复现加 Dreamview 在环，实车还需感知与控制链路配合，展望部分已列出该项。

\section{综合实验\quad{}横向对比}
\textbf{时间}\quad{}8:40 -- 9:30（累计 09:30）\\
\textbf{要点}
\begin{itemize}
  \item 这一页是全文最重量级的数据，请各位老师看主表。
  \item 压力场景 P1--P4\quad{}Baseline 通过率 \MPressureBaselinePassCount/\MPressureScenarioTotal 均触发降级；MIKU \MPressureMikuPassCount/\MPressureScenarioTotal 全通；平均速度 \MPressureBaselineAvgV$\to$\MPressureMikuAvgV\,m/s，QP 总耗时降 84.1\%。
  \item 可比 C1--C4\quad{}两者都能完成通行，MIKU 把 QP 总耗时从 \MComparableBaselineQpTotal\,ms 降到 \MComparableMikuQpTotal\,ms，降 65.6\%。
  \item 汇总 \MScenarioTotal\,场景\quad{}通过 \MSummaryBaselinePassCount/\MScenarioTotal$\to$\MSummaryMikuPassCount/\MScenarioTotal，整体平均速度提升 \MSummarySummaryVImprovement\%。
\end{itemize}
\textbf{过渡}\quad{}消融结果支持三点改进各有贡献。\\
\textbf{可能追问}\quad{}4 个场景是不是偏少。应对，压力加可比共 \MScenarioTotal\,组，覆盖行人、违停车、路锥、施工四类典型威胁物，未来工作会加多车交互场景。

\section{综合实验\quad{}消融结论}
\textbf{时间}\quad{}9:30 -- 9:50（累计 09:50，可留 10 秒缓冲）\\
\textbf{要点}
\begin{itemize}
  \item 六组消融变体是减法消融 leave-one-out 口径，从 MIKU 全量依次关闭单个组件（本节的 C1--C5 指算法组件，与前面可比场景 C1--C4 不同）。
  \item M0 Apollo Baseline\quad{}\AblationScoreMZero\,分，通过 \AblationMZeroPassCount/\AblationMZeroScenarioTotal。
  \item M1 MIKU w/o 组件C1（无 $\tau(s)$ 投影）\quad{}\AblationScoreMOne\,分，\AblationMOnePassCount/\AblationMOneScenarioTotal；动态行人在路径阶段被关成静态快照，P1/P2 均失败。
  \item M2 MIKU w/o 组件C2--C3（无组级决策）\quad{}\AblationScoreMTwo\,分，\AblationMTwoPassCount/\AblationMTwoScenarioTotal；P4 密集施工逐障碍物贪心与水马方向对冲，$s_{end}$ 仅 \AblationMTwoScFourSEndPct\% 即停车。
  \item M3 MIKU w/o 组件C4（统一 $\delta$）\quad{}\AblationScoreMThree\,分，\AblationMThreePassCount/\AblationMThreeScenarioTotal；P3 窄路锥桶统一缓冲压死通行带。
  \item M4 MIKU w/o 组件C5（无走廊注入）\quad{}\AblationScoreMFour\,分，\AblationMFourPassCount/\AblationMFourScenarioTotal；C5 接口在当前压力场景未激活，与 M5 仅差 $\AblationMarginMFourMFive$\,分。
  \item M5 MIKU 全量\quad{}\AblationScoreMFive\,分，\AblationMFivePassCount/\AblationMFiveScenarioTotal。
  \item 结论是路径阶段组件 C1/C2--C3/C4 三组互补，任缺其一即在 P1--P4 中某一场景被压到 \AblationFailedScoreCeil\,分以下；组件 C5 走廊注入的定量贡献需要多动态冲突点场景补测。
\end{itemize}
\textbf{过渡}\quad{}下面看耗时。\\
\textbf{可能追问}\quad{}M4 为什么反而比 M5 的 QP 耗时更低。应对，M5 比 M4 多做一次走廊到速度 QP 的投影，压力场景下该投影尚未起作用但已付出约 \MSummarySummaryExtraMs\,ms 额外开销；一旦遇到多动态冲突点 C5 才会通过收紧速度上界回收耗时。

\section{综合实验\quad{}QP 求解耗时}
\textbf{时间}\quad{}9:50 -- 10:00（累计 10:00）\\
\textbf{要点}
\begin{itemize}
  \item P2 行人加停车差距最大，Baseline QP 总耗时 \MTwoBaselineQpTotal\,ms，MIKU 只要 \MTwoMikuQpTotal\,ms，加速约 20 倍。
  \item C2 行人加停车可比\quad{}\MSixBaselineQpTotal$\to$\MSixMikuQpTotal\,ms，加速约 14 倍。
  \item 压力组均 QP 总耗时 \MPressureBaselineQpTotal$\to$\MPressureMikuQpTotal\,ms，降 84.1\%。
  \item 可比组均 QP 总耗时 \MComparableBaselineQpTotal$\to$\MComparableMikuQpTotal\,ms，降 65.6\%。
  \item 整体 MIKU 比 Baseline 仅多付 \MSummarySummaryExtraMs\,ms 的计算开销，换来 \MPressureMikuPassCount/\MPressureScenarioTotal 个压力场景的突破。
  \item 压力场景 Baseline 因路径阻塞进入 FallbackPath 降级，数值不构成效率比较依据；可比场景的差异主要来自 MIKU 路径已绕开障碍，速度 QP 一次收敛。
\end{itemize}
\textbf{过渡}\quad{}下面是局限性与展望。\\
\textbf{可能追问}\quad{}耗时对比公平吗。应对，论文小结里已列 Baseline 降级模式的客观说明，实验章 chapter 8 单列了耗时复验小节。

\section{研究局限性}
\textbf{时间}\quad{}10:00 -- 10:30（缓冲段）\\
\textbf{要点}
\begin{itemize}
  \item 预测耦合单向，$\tau(s)$ 由上一帧速度反推，没做自车候选回传周车预测。
  \item 纵向重叠组粒度，组内截面求解对链式重叠没做细粒度拆分。
  \item 验证边界，限于 Python 复现加 Dreamview 在环，未覆盖多车交互、恶劣天气与车载嵌入式算力外推。
  \item 鲁棒性边界已用蒙特卡洛量化\quad{}威胁度权重 $\pm\SensPerturbPct\%$ 扰动 \SensTrialCount\,次，P1 到 P3 翻转率 \SensOneMcFlipRate\%，P4 密集施工翻转率 \SensFourMcFlipRate\%，最坏相对偏差 \SensWorstRelDev\%。P4 翻转主要出现在单因子扰动下的 \SensFourSingleFlipCnt\,例决策切换，动态走廊兜底保证可行。
\end{itemize}
\textbf{过渡}\quad{}局限对应未来工作。\\
\textbf{可能追问}\quad{}P4 翻转率为什么偏高。应对，P4 密集施工场景在威胁度评分边界附近的障碍物较多，扰动会让个别障碍物从低威胁翻到中威胁，但动态走廊兜底保证最终可行，最坏偏差控制在 \SensWorstRelDev\%。

\section{研究展望}
\textbf{时间}\quad{}缓冲段\\
\textbf{要点}
\begin{itemize}
  \item \textbf{交互式预测}——把自车候选轨迹回传预测模块，降低预测与规划的单向耦合偏差。
  \item \textbf{路径速度迭代}——将单次前向 $\tau(s)$ 扩展为迭代求解，给出收敛界与单帧预算。
  \item \textbf{组内细粒度}——对纵向链式重叠做细粒度拆分；多间隙宽度接近时引入学习辅助评分。
  \item \textbf{实车部署}——在真实感知与控制链路下评估噪声、延迟与嵌入式算力映射，给出降级与失效安全方案。
\end{itemize>
\textbf{过渡}\quad{}收尾回到结论。\\
\textbf{可能追问}\quad{}为什么要引入学习辅助评分。应对，多个间隙宽度接近时最大间隙策略的选择对外部驾驶习惯不敏感，学习评分可以补充这些因素但不作为决策主干。

\section{主要结论}
\textbf{时间}\quad{}缓冲段\\
\textbf{要点}
\begin{itemize}
  \item 核心结论\quad{}可通过场景下 MIKU 不劣于 Apollo Baseline，轨迹平顺性与 QP 求解开销均有改善；Baseline 触发 BLOCKED 的压力场景下，MIKU 通过差异化 $\delta_i$、组级 $\mathbf{d}^*$、时空走廊 $\tau(s)$ 三条改进把通行率从 \MPressureBaselinePassCount/\MPressureScenarioTotal 提升至 \MPressureMikuPassCount/\MPressureScenarioTotal。
  \item 方法\quad{}三条改进共用同一组 $(u_i, v_i, \delta, \tau)$ 符号，源码层只动 PathBoundsDecider 一步，上下游 PJPO、PJSO 与 Control 链路零改动。
  \item 效率\quad{}整体平均速度提升 \MSummarySummaryVImprovement\%，压力平均速度 \MPressureBaselineAvgV$\to$\MPressureMikuAvgV\,m/s。
  \item 计算\quad{}压力 QP 耗时 \MPressureBaselineQpTotal$\to$\MPressureMikuQpTotal\,ms（降 84.1\%），可比 \MComparableBaselineQpTotal$\to$\MComparableMikuQpTotal\,ms（降 65.6\%），整体仅多付 \MSummarySummaryExtraMs\,ms。
  \item 消融\quad{}减法消融下，关闭路径阶段任一组件（M1/M2/M3）即在 P1--P4 中某场景失败，评分被压到 \AblationFailedScoreCeil\,分以下；M4（w/o 组件 C5）\AblationScoreMFour\,分，与 M5 仅差 $\AblationMarginMFourMFive$\,分，组件 C5 的边际贡献待多冲突点场景补测，路径阶段组件 C1/C2--C3/C4 互补。
  \item 鲁棒\quad{}威胁度权重 $\pm\SensPerturbPct\%$ 扰动 \SensTrialCount\,次蒙特卡洛，最坏相对偏差 \SensWorstRelDev\%。
  \item 已完成 C++ 增量接入。
\end{itemize}
\textbf{过渡}\quad{}最后致谢。\\
\textbf{可能追问}\quad{}下一步投稿吗。应对，相关工作会整理成英文稿件投 IV 或 ITSC 系列会议。

\section{致谢}
\textbf{时间}\quad{}10:20 -- 10:30（累计 10:30，留 30 秒缓冲）\\
\textbf{要点}
\begin{itemize}
  \item 感谢孙成娇老师四年来指导。
  \item 感谢计算机工程学院老师们在答辩准备上给的意见。
  \item 答辩到此结束，恳请各位老师批评指正。
\end{itemize}

\end{document}
